\chapter{Chebyshev's Bias}\label{ch:bias}
\begin{center}\itshape
In the race between the primes $4n+1$ and $4n+3$, the squares do not run.\\
That is enough to make one side lead $99.59\%$ of the time,\\
and not enough to make it lead for ever.
\end{center}

In 1853, Pafnuty Chebyshev sent a short letter to Paul Fuss, the permanent secretary of the Saint Petersburg Academy. Its point, in substance, was that there are noticeably more primes of the form $4n+3$ than of the form $4n+1$, and that the excess can be seen in a weighted sum which, Chebyshev asserted, tends to $-\infty$. The observation was right, and puzzling: Dirichlet had shown in 1837 that both forms contain infinitely many primes, and by the prime number theorem for progressions (proved in 1896) they contain asymptotically the same number. Tome~5, Chapter~4 drew the race. Up to $26860$ the form $4n+3$ is ahead or tied at every single $x$. John Leech found the first lead change in 1957, with one of the first computers: at $x=26861$ the form $4n+1$ goes ahead by one prime, and loses the lead at once. We ran the race further for this chapter. Up to $10^9$ there are $25\,423\,491$ primes $4n+1$ and $25\,424\,042$ primes $4n+3$, and the form $4n+1$ has led for only $0.04\%$ of the way, measured on the logarithmic scale that suits the problem.

Two things needed explaining: why the bias, and why it is not total. Tome~5, Chapter~7 answered the first question with an identity and a heuristic. The race between the primes is the race between the prime \emph{powers}, which is fair, plus a head start for $4n+3$ coming from the squares of primes, all of which are $\equiv1\pmod4$ and are counted on the wrong side. In a fair race, an oscillating sum over the zeros of $L(s,\chi_{-4})$ decides the leader; the head start is of exactly the size of that oscillation, so it biases the race without settling it.

The second question was answered first. In 1914 John Edensor Littlewood proved that the lead changes infinitely often, and that the same holds for the more famous race between $\pi(x)$ and $\operatorname{li}(x)$, in which every computation then available had $\operatorname{li}(x)$ ahead. In 1918 G.~H.~Hardy and Littlewood, and independently Edmund Landau, showed that Chebyshev's weighted sum does tend to $-\infty$ if and only if the generalized Riemann hypothesis holds for $L(s,\chi_{-4})$: Chebyshev's ``observation'' was a deep conjecture. And in 1994 Michael Rubinstein and Peter Sarnak turned the heuristic into a theorem. Assuming the generalized Riemann hypothesis and a hypothesis of linear independence for the zeros, they proved that the race has a limiting distribution, the law of a sum of independent random waves, and they computed the share of the time that $4n+3$ leads: $99.59\%$.

This chapter proves their theorem. It is a fitting continuation of Chapter~\ref{ch:bohrjessen}, with the zeros in the role of the primes: there, $\log\zeta(\sigma+it)$ was a sum of waves $p^{-it}$ with independent phases; here, the error term in the race is a sum of waves $x^{i\gamma}$, and the hypothesis that the ordinates $\gamma$ are linearly independent plays the part of unique factorisation.

\begin{reminder}[from Tomes~5, 6, 12 and~13]
\emph{The identity} (Tome~5, Proposition~7.11). With $\chi=\chi_{-4}$ and $\theta(x;4,a)=\sum_{p\le x,\,p\equiv a\,(4)}\log p$, for $x\ge2$,
$\theta(x;4,3)-\theta(x;4,1)=\theta(\sqrt x)-\psi(x,\chi)+O(x^{1/3}\log x)$, where $\psi(x,\chi)=\sum_{n\le x}\chi(n)\Lambda(n)$.
\emph{Counting zeros locally} (Tome~5, Lemma~6.6). Let $f$ be holomorphic on a neighbourhood of $|z|\le R$, with $f(0)\ne0$ and $|f(z)|\le M|f(0)|$ for $|z|\le R$, $M\ge2$, and let $\mathcal Z$ be its zeros in $|z|\le\frac R2$. Then $\#\mathcal Z\le\frac{\log M}{\log2}$, and $|\frac{f'}f(z)-\sum_{\rho\in\mathcal Z}\frac1{z-\rho}|\le\frac{17\log M}R$ for $|z|\le\frac R4$.
\emph{The functional equation} (Tome~12, Corollary~9.21). For $\chi=\chi_{-4}$, which is primitive and odd, $\xi(s)=(\frac4\pi)^{s/2}\Gamma(\frac{s+1}2)L(s,\chi)$ is entire and $\xi(s)=\xi(1-s)$. (The root number is $\frac{\tau(\chi)}{i\sqrt4}=\frac{2i}{2i}=1$.)
\emph{Stirling's formula} (Tome~5, Theorem~6.1), in the form: for $\sigma$ in a fixed compact interval and $|t|\ge1$, $|\Gamma(\sigma+it)|\asymp|t|^{\sigma-1/2}e^{-\pi|t|/2}$, and $\frac{\Gamma'}\Gamma(w)=O(\log(|w|+2))$ for $\operatorname{Re}w\ge\frac12$.
\emph{Perron's formula for characters} (Tome~13, Lemma~2.10). For $x\ge3$, $T\ge2$ and $a=1+\frac1{\log x}$, $\psi(x,\chi)=\frac1{2\pi i}\int_{a-iT}^{a+iT}(-\frac{L'}L(s,\chi))\frac{x^s}s\dd s+O(\frac{x\log^2x}T+\log x)$.
\emph{Weak convergence} (Tome~6, Theorems~11.4 and~11.7). $\mu_n\Rightarrow\mu$ if and only if $\int f\dd\mu_n\to\int f\dd\mu$ for every bounded Lipschitz $f$, if and only if $\mu_n(A)\to\mu(A)$ for every Borel $A$ with $\mu(\partial A)=0$; and then $\liminf\mu_n(U)\ge\mu(U)$ for open $U$. A tight sequence of probabilities on $\R$ has a weakly convergent subsequence (Prokhorov).
\end{reminder}

Throughout the chapter $\chi=\chi_{-4}$ is the non-principal character modulo $4$, $L(s)=L(s,\chi)$, and $\rho=\beta+i\gamma$ runs over the \emph{non-trivial zeros} of $L$, those with $0<\beta<1$, counted with multiplicity. Since $L$ has an Euler product for $\sigma>1$ and does not vanish on $\sigma=1$ (Tome~13, Theorem~2.3), the functional equation shows that its other zeros are the \emph{trivial zeros} $-1,-3,-5,\dots$, the poles of $\Gamma(\frac{s+1}2)$, and that $L(it)\ne0$: indeed $\xi(s)=\xi(1-s)\ne0$ for $\sigma\le0$. Since $\chi$ is real, $L(\overline s)=\overline{L(s)}$, so the zeros are symmetric about the real axis, with multiplicities.

\section{The explicit formula for \texorpdfstring{$L(s,\chi_{-4})$}{L(s,chi-4)}}

Tome~5 proved the explicit formula for $\zeta$ and stated its analogue for $L(s,\chi_{-4})$ without proof. We need the analogue, so we prove it. Hadamard's product is not needed: local information on the zeros suffices.

\begin{theorem}\label{thm:explicitchi}
\leavevmode
\begin{enumerate}[label=(\alph*)]
\item For every real $t$, the number of zeros $\rho$ with $|\gamma-t|\le1$ is $O(\log(|t|+2))$.
\item Uniformly for $-1\le\sigma\le2$ and $|t|\ge2$, $\dfrac{L'}L(\sigma+it)=\displaystyle\sum_{|\gamma-t|\le1}\frac1{s-\rho}+O(\log|t|)$.
\item For $x\ge2$ and $T\ge3$, with $\psi_0(x,\chi)$ the average of the left and right limits of $\psi(\cdot,\chi)$ at $x$ and $\langle x\rangle$ the distance from $x$ to the nearest prime power other than $x$,
\[
\psi_0(x,\chi)=-\sum_{|\gamma|<T}\frac{x^\rho}\rho-\frac{L'}L(0)+\sum_{k\ge1}\frac{x^{1-2k}}{2k-1}+O\Bigl(\frac{x\log^2(xT)}T+\log x\min\Bigl(1,\frac x{T\langle x\rangle}\Bigr)\Bigr).
\]
\end{enumerate}
The implied constants are absolute.
\end{theorem}
\begin{proof}
\emph{Growth.} Let $S(u)=\sum_{n\le u}\chi(n)\in\{0,1\}$. By partial summation, $L(s)=s\int_1^\infty S(u)u^{-s-1}\dd u$ for $\sigma>0$, so $|L(s)|\le\frac{|s|}\sigma\le2|s|$ for $\sigma\ge\frac12$. For $\sigma<\frac12$, the functional equation gives
\[
L(s)=\Bigl(\frac4\pi\Bigr)^{\frac12-s}\frac{\Gamma(\frac{2-s}2)}{\Gamma(\frac{s+1}2)}L(1-s),
\]
and by Stirling's formula the ratio of Gamma functions is $O((|t|+2)^{\frac12-\sigma})$ for $-11\le\sigma\le\frac12$ and $|t|\ge1$; for $|t|\le1$ it is bounded on that strip, because $\Gamma(\frac{2-s}2)$ has no pole there and $\frac1\Gamma$ is entire. Hence $|L(s)|\le C(|t|+2)^{13}$ for $-11\le\sigma\le15$.

\emph{(a) and (b).} Fix $t$ and let $f(z)=L(z+2+it)$, $R=12$. The Euler product gives $|f(0)|=|L(2+it)|\ge\prod_p(1+p^{-2})^{-1}=\frac{\zeta(4)}{\zeta(2)}>\frac12$, so the hypothesis of Tome~5, Lemma~6.6 holds with $\log M=O(\log(|t|+2))$. Every zero with $|\gamma-t|\le1$ lies within $\sqrt{2^2+1}<6=\frac R2$ of $2+it$, which proves (a). For (b), let $|t|\ge2$ and $s=\sigma+it$ with $-1\le\sigma\le2$, so $|s-(2+it)|\le3=\frac R4$. The lemma gives $\frac{L'}L(s)=\sum_{\rho\in\mathcal Z}\frac1{s-\rho}+O(\log|t|)$, where $\mathcal Z$ is the set of zeros (trivial ones included) within $6$ of $2+it$. There are $O(\log|t|)$ of them by the same lemma, and those with $|\gamma-t|>1$, or trivial (then $|s-\rho|\ge|t|\ge2$), contribute at most $1$ each.

\emph{Good heights.} Let $T\ge3$. By (a) there are at most $C\log T$ ordinates in $[T,T+1]$, so some $T'\in[T,T+1]$ is at distance at least $\frac1{2(C\log T+1)}$ from all of them. By (b), for $-1\le\sigma\le2$, $|\frac{L'}L(\sigma\pm iT')|\le\sum_{|\gamma\mp T'|\le1}\frac1{|T'\mp\gamma|}+O(\log T)=O(\log^2T)$, using (a) again and the symmetry of the zeros. For $\sigma\ge2$, $|\frac{L'}L(s)|\le-\frac{\zeta'}\zeta(2)$.

\emph{The left half-plane.} Taking logarithmic derivatives in the functional equation,
\[
\frac{L'}L(s)=-\log\frac4\pi-\frac12\frac{\Gamma'}\Gamma\Bigl(\frac{2-s}2\Bigr)-\frac12\frac{\Gamma'}\Gamma\Bigl(\frac{s+1}2\Bigr)-\frac{L'}L(1-s).
\]
Let $\sigma\le-1$. Then $|\frac{L'}L(1-s)|\le-\frac{\zeta'}\zeta(2)$, and $\frac{\Gamma'}\Gamma(\frac{2-s}2)=O(\log(|s|+2))$. With $w=\frac{s+1}2$, the reflection formula $\frac{\Gamma'}\Gamma(w)=\frac{\Gamma'}\Gamma(1-w)-\pi\cot\pi w$ and $\operatorname{Re}(1-w)\ge1$ give $\frac{\Gamma'}\Gamma(w)=O(\log(|s|+2))$ provided $s$ stays at distance at least $1$ from the odd integers, where $\cot\pi w$ is bounded. So $\frac{L'}L(s)=O(\log(|s|+2))$ on the lines $\sigma=-2K$ ($K\ge1$ an integer) and on $\{\sigma\le-1,\ |t|\ge3\}$.

\emph{The rectangle.} Tome~5, Proposition~7.2 (the truncated Perron formula) holds for $\psi_0(x,\chi)$, with $-\frac{L'}L$ in place of $-\frac{\zeta'}\zeta$ and the same error term, because its proof uses only $|\chi(n)\Lambda(n)|\le\Lambda(n)$; this is how Tome~13, Lemma~2.10 is proved. Let $a=1+\frac1{\log x}$ and $F(s)=-\frac{L'}L(s)\frac{x^s}s$. Apply the residue theorem to the rectangle with vertices $a\pm iT'$, $-2K\pm iT'$. Inside, $F$ has poles at the non-trivial zeros, with residue $-m\frac{x^\rho}\rho$ at a zero of multiplicity $m$; at $s=0$, with residue $-\frac{L'}L(0)$ (recall $L(0)\ne0$); and at the trivial zeros $1-2k$, $1\le k\le K$, which are simple (simple poles of $\Gamma$, and $\xi(1-2k)=\xi(2k)\ne0$), with residue $-\frac{x^{1-2k}}{1-2k}$. The horizontal sides contribute $O(\frac{x\log^2T}{T\log x})$ exactly as in Tome~5, Theorem~7.5: on $-1\le\sigma\le a$ the integrand is $O(\frac{\log^2T}Tx^\sigma)$, and on $\sigma\le-1$ it is $O(\frac{\log(|\sigma|+T)}Tx^\sigma)$. The left side is $O(T\log(K+T)\frac{x^{-2K}}{K})$, which tends to $0$ as $K\to\infty$. Combining with Perron's formula at height $T'$ gives (c) with $T'$ in place of $T$, the errors adding up to $O(\frac{x\log^2x}T+\log x\min(1,\frac x{T\langle x\rangle})+\frac{x\log^2T}{T\log x})$. Finally, the zeros with $T\le|\gamma|<T'$ are $O(\log T)$ in number by (a), and each term is at most $\frac xT$ in modulus.
\end{proof}

\begin{corollary}\label{cor:grhchi}
Assume the generalized Riemann hypothesis for $\chi$ (GRH): every non-trivial zero has $\beta=\frac12$. Then
\begin{enumerate}[label=(\alph*)]
\item $\psi(x,\chi)=O(\sqrt x\log^2x)$ for $x\ge3$;
\item $\sum_\rho\frac1{|\rho|}=\infty$ (this part does not need GRH);
\item for $T_0\ge0$, $\displaystyle\sum_{|\gamma|\ge T_0}\frac{\log^2(|\gamma|+2)}{(1+|\gamma|)^2}=O\Bigl(\frac{\log^3(T_0+2)}{T_0+1}\Bigr)$; in particular $\sum_\rho|\rho|^{-2}<\infty$.
\end{enumerate}
\end{corollary}
\begin{proof}
(c) Group the zeros by $n\le|\gamma|<n+1$; by Theorem~\ref{thm:explicitchi}(a) there are $O(\log(n+2))$ in each group, so the sum is $O(\sum_{n\ge T_0-1}\frac{\log^3(n+2)}{(n+1)^2})$, which is $O(\frac{\log^3(T_0+2)}{T_0+1})$ by comparison with an integral. Since $|\rho|\ge\max(\frac12,|\gamma|)\ge\frac{1+|\gamma|}3$, the last claim follows with $T_0=0$. (a) Take $T=x$ in Theorem~\ref{thm:explicitchi}(c): $|x^\rho|=\sqrt x$ and, grouping in the same way, $\sum_{|\gamma|<x}\frac1{|\rho|}=O(\sum_{n<x}\frac{\log(n+2)}{n+1})=O(\log^2x)$; also $|\psi-\psi_0|\le\log x$. (b) Suppose $\sum_\rho\frac1{|\rho|}<\infty$. For $1<x<x'$ in a compact interval, $|x^\rho|\le x'$, so the series $\sum_\rho\frac{x^\rho}\rho$ converges uniformly there and defines a continuous function of $x>1$. Letting $T\to\infty$ in Theorem~\ref{thm:explicitchi}(c) for $x$ not a prime power (then $\langle x\rangle>0$ and the error tends to $0$), $\psi_0(x,\chi)$ equals a continuous function of $x$ on the set of such $x$. But $\psi(\cdot,\chi)$ jumps by $\chi(5)\log5=\log5$ at $x=5$ and is constant on $(4,5)$ and on $(5,7)$: a contradiction.
\end{proof}

\section{The race as a sum of waves}

From now on we assume GRH for $\chi$, and write each zero as $\rho=\frac12+i\gamma$. The object of study is the normalised lead of the form $4n+3$, in the variable $u=\log x$:
\[
E(u)=\frac u{e^{u/2}}\bigl(\pi(e^u;4,3)-\pi(e^u;4,1)\bigr)\qquad(u\ge\log2).
\]
Its model, for $T\ge3$, is the trigonometric polynomial
\[
E_T(u)=1+\sum_{|\gamma|<T}\frac{e^{i\gamma u}}\rho=1+2\sum_{0<\gamma<T}\operatorname{Re}\frac{e^{i\gamma u}}\rho+\sum_{\gamma=0}\frac1\rho,
\]
which is real because the zeros are symmetric. (The last sum is empty: $L(\frac12)=0.6677\ldots\ne0$, but we shall not need this.) The plan is to show that $E$ is close to $E_T$ in mean square, with an error that tends to $0$ as $T\to\infty$, and the first step is a mean value estimate for sums over zeros.

\begin{lemma}\label{lem:zerowaves}
Assume GRH. There is an absolute constant $C$ such that for $0\le T_0<T$ and every real $V$,
\[
\int_V^{V+1}\Bigl|\sum_{T_0\le|\gamma|<T}\frac{e^{i\gamma u}}\rho\Bigr|^2\dd u\le C\frac{\log^3(T_0+2)}{T_0+1}.
\]
\end{lemma}
\begin{proof}
Expanding the square and integrating term by term, the left side is at most $\sum_{\gamma,\gamma'}\frac1{|\rho||\rho'|}\min(1,\frac2{|\gamma-\gamma'|})$, both variables running over $T_0\le|\gamma|<T$. We have $\min(1,\frac2{|d|})\le\frac4{1+|d|}$ and $\frac1{|\rho'|}\le\frac3{1+|\gamma'|}$, so for fixed $\gamma$ the inner sum is at most $12\,s(\gamma)$, where
\[
s(\gamma)=\sum_{\gamma'}\frac1{(1+|\gamma-\gamma'|)(1+|\gamma'|)}.
\]
Let $g=|\gamma|$. Since $(1+|\gamma-\gamma'|)+(1+|\gamma'|)\ge2+g$, each term is at most $\frac1{2+g}(\frac1{1+|\gamma-\gamma'|}+\frac1{1+|\gamma'|})$. Over the zeros with $|\gamma'|\le2g+2$, grouped in unit intervals of $O(\log(g+2))$ zeros each, the sums of $\frac1{1+|\gamma-\gamma'|}$ and of $\frac1{1+|\gamma'|}$ are $O(\log^2(g+2))$. Over the zeros with $|\gamma'|>2g+2$, $|\gamma-\gamma'|\ge\frac{|\gamma'|}2$ and each term is at most $\frac2{(1+|\gamma'|)^2}$; these sum to $O(\frac{\log(g+2)}{g+1})$. Hence $s(\gamma)=O(\frac{\log^2(g+2)}{g+1})$, and the left side is
\[
O\Bigl(\sum_{|\gamma|\ge T_0}\frac{\log^2(|\gamma|+2)}{(1+|\gamma|)^2}\Bigr)=O\Bigl(\frac{\log^3(T_0+2)}{T_0+1}\Bigr)
\]
by Corollary~\ref{cor:grhchi}(c).
\end{proof}

The case $T_0=0$ says that the waves have a bounded mean square on every window of length $1$, however many of them we take. This is not obvious: the series $\sum_\rho\frac{e^{i\gamma u}}\rho$ does not converge absolutely, because $\sum_\rho\frac1{|\rho|}$ diverges (Corollary~\ref{cor:grhchi}(b)).

\begin{proposition}\label{prop:racel2}
Assume GRH. There is an absolute constant $C$ such that for $T\ge3$ and $U\ge2+\log T$,
\[
\frac1U\int_2^U|E(u)-E_T(u)|^2\dd u\le C\Bigl(\frac{\log^3T}T+\frac{\log T}U\Bigr).
\]
\end{proposition}
\begin{proof}
\emph{Weighted primes.} Let
\[
F(u)=e^{-u/2}\bigl(\theta(e^u;4,3)-\theta(e^u;4,1)\bigr).
\]
Fix an integer $V\ge2$ and $u\in[V,V+1]$, $x=e^u$. Theorem~\ref{thm:explicitchi}(c) with $T_V=e^{2V+2}\ge x^2$, together with $|\psi-\psi_0|\le\log x$, gives
\[
\psi(x,\chi)=-\sum_{|\gamma|<T_V}\frac{x^\rho}\rho+O(V).
\] The prime number theorem in the form $\theta(y)=y+O(\frac y{\log^2y})$ (Tome~5, Theorem~7.8) gives $\theta(\sqrt x)=\sqrt x+O(\frac{\sqrt x}{u^2})$. With the identity of the reminder,
\[
F(u)=1+\sum_{|\gamma|<T_V}\frac{e^{i\gamma u}}\rho+O(u^{-2}).
\]
So $F-E_T$ is $\pm$ a sum over $\min(T,T_V)\le|\gamma|<\max(T,T_V)$, plus $O(V^{-2})$. By Lemma~\ref{lem:zerowaves}, $\int_V^{V+1}|F-E_T|^2\le C(\frac{\log^3T}T+V^{-4})$ when $T_V\ge T$, and $\int_V^{V+1}|F-E_T|^2\le C$ in all cases. The second bound is used for the $O(\log T)$ values of $V$ with $T_V<T$. Summing over $V<U$,
\[
\int_2^U|F-E_T|^2\dd u\le C\Bigl(\frac{U\log^3T}T+\log T+1\Bigr).
\]
The same argument with Lemma~\ref{lem:zerowaves} for $T_0=0$ shows $\int_V^{V+1}|F|^2\le C$ for all $V\ge2$, and $F$ is bounded on $[\log2,2]$.

\emph{From $\theta$ to $\pi$.} Let $D(x)=\theta(x;4,3)-\theta(x;4,1)$. By partial summation, $\pi(x;4,a)=\frac{\theta(x;4,a)}{\log x}+\int_2^x\frac{\theta(t;4,a)}{t\log^2t}\dd t$, so with $t=e^v$,
\[
E(u)=F(u)+G(u),\qquad G(u)=ue^{-u/2}\int_{\log2}^u\frac{e^{v/2}F(v)}{v^2}\dd v.
\]
By Cauchy--Schwarz, with the weight $e^{v/2}v^{-2}$, and $\int_{\log2}^ue^{v/2}v^{-2}\dd v\le C\frac{e^{u/2}}{u^2}$,
\[
|G(u)|^2\le C e^{-u/2}\int_{\log2}^u\frac{e^{v/2}|F(v)|^2}{v^2}\dd v.
\]
Integrating over $u\in[2,U]$ and exchanging the integrals, $\int_2^U|G|^2\le C\int_{\log2}^U|F(v)|^2v^{-2}\dd v\le C'$, by the window bounds for $F$. Since $|E-E_T|^2\le2|F-E_T|^2+2|G|^2$, the proposition follows.
\end{proof}

So, in mean square, the normalised race is $1+\sum_\gamma\frac{e^{i\gamma u}}\rho$: a constant head start of exactly $1$ for the form $4n+3$, the squares' contribution, plus an oscillation of mean zero. Note how the conversion from $\theta$ to $\pi$ costs nothing at this scale: the term $G$ is smaller than $F$ by a factor $u$, in mean square.

\section{The limiting distribution}

\begin{lemma}[Kronecker--Weyl for real frequencies]\label{lem:kwreal}
Let $\lambda_1,\dots,\lambda_n$ be real. For every continuous $\Phi$ on the torus $\mathbb T^n$, the average $\frac1U\int_0^U\Phi(e^{i\lambda_1u},\dots,e^{i\lambda_nu})\dd u$ has a limit as $U\to\infty$. If $\lambda_1,\dots,\lambda_n$ are linearly independent over $\Q$, the limit is $\mathbb E\Phi(Z_1,\dots,Z_n)$, with $Z_1,\dots,Z_n$ independent and uniform on the unit circle.
\end{lemma}
\begin{proof}
For a monomial $\Phi(z)=z_1^{m_1}\cdots z_n^{m_n}$ the average is $\frac1U\int_0^Ue^{i(m\cdot\lambda)u}\dd u$, which is $1$ if $m\cdot\lambda=0$ and at most $\frac2{U|m\cdot\lambda|}$ in modulus otherwise. So the limit exists for trigonometric polynomials; under linear independence $m\cdot\lambda=0$ only for $m=0$, and the limit is $\mathbb E\Phi(Z)$, which is also $1$ or $0$. Trigonometric polynomials are uniformly dense in $C(\mathbb T^n)$ (as in Lemma~\ref{lem:kw}), and if $|\Phi-\Phi_k|\le\varepsilon$ uniformly, the averages of $\Phi$ and $\Phi_k$ differ by at most $\varepsilon$ for every $U$; so the averages of $\Phi$ form a Cauchy family.
\end{proof}

\begin{theorem}[Wintner, 1935; Rubinstein and Sarnak, 1994]\label{thm:racelimit}
Assume GRH for $\chi_{-4}$. There is a probability measure $\mu$ on $\R$ such that the law of $E(u)$, for $u$ uniform in $[2,U]$, converges weakly to $\mu$ as $U\to\infty$: for every bounded continuous $f$,
\[
\frac1{U-2}\int_2^Uf(E(u))\dd u\longrightarrow\int f\dd\mu.
\]
\end{theorem}
\begin{proof}
Let $P_U$ be the law of $E(u)$ and $P_U^T$ that of $E_T(u)$, for $u$ uniform in $[2,U]$. For bounded continuous $f$, $f(E_T(u))=\Phi(e^{i\gamma u})_{|\gamma|<T}$ with $\Phi$ continuous on a torus, so by Lemma~\ref{lem:kwreal} (shifting $u$ by $2$ does not change the limit) $\int f\dd P_U^T$ has a limit $\ell_T(f)$ as $U\to\infty$. By Lemma~\ref{lem:zerowaves} with $T_0=0$, applied on each window $[V,V+1]$, $\int x^2\dd P_U^T\le C$ for all $U$ and $T$; so the family $(P_U^T)_U$ is tight, some sequence $P_{U_n}^T$ converges weakly to a probability $\mu_T$, and $\int f\dd\mu_T=\ell_T(f)$. Moreover $\int x^2\dd\mu_T\le C$.

Let $f$ be bounded by $1$ and $1$-Lipschitz. Then by Proposition~\ref{prop:racel2} and Cauchy--Schwarz,
\[
\Bigl|\int f\dd P_U-\int f\dd P_U^T\Bigr|\le\frac1{U-2}\int_2^U|E-E_T|\le C\Bigl(\frac{\log^3T}T+\frac{\log T}U\Bigr)^{1/2}.
\]
Letting $U\to\infty$, the $\limsup$ and the $\liminf$ of $\int f\dd P_U$ lie within $C(\frac{\log^3T}T)^{1/2}$ of $\int f\dd\mu_T$, for every $T$. Hence $\int f\dd P_U$ converges to some $\ell(f)$, and $\int f\dd\mu_T\to\ell(f)$ as $T\to\infty$. The family $(\mu_T)$ is tight. By Prokhorov's theorem some sequence $\mu_{T_j}$ converges weakly to a probability $\mu$, and then $\int f\dd\mu=\ell(f)$ for all bounded Lipschitz $f$. So $P_U\Rightarrow\mu$ by the portmanteau theorem.
\end{proof}

Without further hypotheses, little can be said about $\mu$. Its first two moments can be computed, and it is the law of a function on a torus. But if the ordinates of the zeros had rational relations, the waves could conspire, and $\mu$ could even fail to give the race to the form $4n+3$. Rubinstein and Sarnak made the assumption that no such conspiracy exists.

\medskip\noindent\textbf{The linear independence hypothesis (LI).} \emph{The positive ordinates $\gamma>0$ of the zeros of $L(s,\chi_{-4})$, listed with multiplicity, are linearly independent over $\Q$, and $L(\frac12,\chi_{-4})\ne0$.}
\medskip

In particular LI implies that the zeros are simple. Nobody knows how to prove it, or even that the first two ordinates, $6.0209\ldots$ and $10.2437\ldots$, have an irrational ratio; but no relation has ever been found, and LI is the kind of statement one expects of numbers that ``have no reason'' to be related. It plays here exactly the part that unique factorisation played in Chapter~\ref{ch:bohrjessen}, where it made the waves $p^{-it}$ independent.

\begin{theorem}[Rubinstein and Sarnak, 1994]\label{thm:rubinsteinsarnak}
Assume GRH and LI for $\chi_{-4}$. Let $(\Theta_\gamma)_{\gamma>0}$ be independent and uniform on $[0,2\pi)$. Then the series
\[
X=1+\sum_{\gamma>0}\frac{2\cos\Theta_\gamma}{\sqrt{\frac14+\gamma^2}}
\]
converges almost surely and in $L^2$, and $\mu$ is the law of $X$. It has mean $1$ and variance $V=\sum_\rho\frac1{|\rho|^2}$; it is symmetric about $1$; it has no atoms; and every open interval has positive $\mu$-measure.
\end{theorem}
\begin{proof}
\emph{The series.} The terms are independent, centred and bounded, with variances $\frac{4\cdot\frac12}{|\rho|^2}=\frac2{|\rho|^2}$, and $\sum_{\gamma>0}\frac2{|\rho|^2}=\sum_\rho\frac1{|\rho|^2}=V<\infty$ by Corollary~\ref{cor:grhchi}(c). A series of independent centred variables with summable variances converges in $L^2$ (its partial sums are Cauchy, by orthogonality) and almost surely (Tome~10, Theorem~2.5). The mean is $1$ and the variance $V$.

\emph{Identification.} Since $L(\frac12)\ne0$, $E_T(u)=1+\sum_{0<\gamma<T}\frac{2\cos(\gamma u-\arg\rho)}{|\rho|}$. By LI and Lemma~\ref{lem:kwreal}, the limit law $\mu_T$ of $E_T$ is that of
\[
1+\sum_{0<\gamma<T}\frac{2\cos(\Theta_\gamma-\arg\rho)}{|\rho|},
\]
which is the law of the partial sum $X_T$ of the series, because $\Theta_\gamma-\arg\rho$ is again uniform modulo $2\pi$. As $X_T\to X$ in $L^2$, $\mu_T$ converges weakly to the law of $X$. In the proof of Theorem~\ref{thm:racelimit}, $\int f\dd\mu_T\to\int f\dd\mu$ for bounded Lipschitz $f$, so $\mu$ is the law of $X$.

\emph{Symmetry.} Replacing every $\Theta_\gamma$ by $\Theta_\gamma+\pi$ does not change the joint law of the phases and changes $X-1$ into $1-X$.

\emph{No atoms.} Let $\gamma_1$ be the smallest positive ordinate and $Y=X-\frac{2\cos\Theta_{\gamma_1}}{|\rho_1|}$, which is independent of $\Theta_{\gamma_1}$. For every $c$, $P(X=c)=\mathbb E\bigl[P(2\cos\Theta_{\gamma_1}=|\rho_1|(c-Y)\mid Y)\bigr]=0$, since $\cos\Theta_{\gamma_1}$ takes each value with probability $0$.

\emph{Full support.} Let $a\in\R$ and $\varepsilon>0$. Since $\sum_\rho\frac1{|\rho|}=\infty$ (Corollary~\ref{cor:grhchi}(b)), there is $N$ such that $A_N=\sum_{0<\gamma\le\gamma_N}\frac2{|\rho|}>|a-1|$, where $\gamma_1<\gamma_2<\cdots$ are the positive ordinates; take $N$ also so large that $W_N=\sum_{\gamma>\gamma_N}\frac2{|\rho|^2}<\frac{\varepsilon^2}8$. The function $\theta\mapsto1+\sum_{j\le N}\frac{2\cos\theta_j}{|\rho_j|}$ is continuous on the connected torus $[0,2\pi)^N$ and takes the values $1\pm A_N$, so it takes the value $a$, and it is within $\frac\varepsilon2$ of $a$ on a neighbourhood of that point: the head $X_{\gamma_N}$ of the series satisfies $P(|X_{\gamma_N}-a|<\frac\varepsilon2)>0$. By Chebyshev's inequality the tail satisfies $P(|X-X_{\gamma_N}|\ge\frac\varepsilon2)\le\frac{4W_N}{\varepsilon^2}<\frac12$. The head and the tail are independent, so $P(|X-a|<\varepsilon)>0$.
\end{proof}

\begin{corollary}\label{cor:biasdensity}
Assume GRH and LI for $\chi_{-4}$. The logarithmic density
\[
\delta(4;3,1)=\lim_{X\to\infty}\frac1{\log X}\int_2^X\mathbf 1\{\pi(x;4,3)>\pi(x;4,1)\}\frac{\dd x}x
\]
exists, equals $\mu((0,\infty))=P(X>0)$, and satisfies $\frac12<\delta(4;3,1)<1$. In particular $\pi(x;4,1)>\pi(x;4,3)$ for arbitrarily large $x$.
\end{corollary}
\begin{proof}
With $x=e^u$ the integral is $\int_{\log2}^{\log X}\mathbf 1\{E(u)>0\}\dd u$, and the interval $[\log2,2]$ does not matter in the limit. The boundary of $(0,\infty)$ is $\{0\}$, which has $\mu$-measure $0$, so Theorem~\ref{thm:racelimit} and the portmanteau theorem give the limit $\mu((0,\infty))$. By symmetry and the absence of atoms, $P(X\le0)=P(X-1\le-1)=P(X-1\ge1)=\frac12P(|X-1|\ge1)$, which is less than $\frac12$ because $P(|X-1|<1)>0$ (full support); so $\delta>\frac12$. And $\delta<1$ because $\mu((-\infty,0))>0$. Finally, the open set $(-\infty,0)$ has $\liminf_U P_U((-\infty,0))\ge\mu((-\infty,0))>0$: the set of $x$ where $4n+1$ leads has positive lower logarithmic density, so it is unbounded.
\end{proof}

How large is $\delta(4;3,1)$? Chebyshev's inequality already gives $\delta\ge1-\frac V2$ (Exercise~\ref{ex:chebbound}), and $V$ can be computed exactly from $L$ alone: $V=\log\frac4\pi-\gamma+2\frac{L'}L(1,\chi)=0.1556\ldots$ (Exercise~\ref{ex:variance}), so $\delta\ge0.92$. The exact value needs the zeros. The characteristic function of $X$ is a product,
\[
\mathbb Ee^{i\xi X}=e^{i\xi}\prod_{\gamma>0}J_0\Bigl(\frac{2\xi}{|\rho|}\Bigr),\qquad J_0(z)=\frac1{2\pi}\int_0^{2\pi}e^{iz\cos\theta}\dd\theta,
\]
and since $X-1$ is symmetric (Exercise~\ref{ex:gilpelaez}),
\[
\delta=\frac12+\frac1\pi\int_0^\infty\frac{\sin\xi}\xi\prod_{\gamma>0}J_0\Bigl(\frac{2\xi}{|\rho|}\Bigr)\dd\xi.
\] With a few thousand zeros, Rubinstein and Sarnak found
\[
\delta(4;3,1)=0.9959\ldots
\]
The notebook of this chapter repeats the computation with the $3130$ zeros below height $3000$ and a Gaussian model for the rest, and finds $0.99593$; Figure~\ref{fig:biasdensity} shows the density of $\mu$.

\begin{figure}[ht]
\centering\fig{ch09-race}
\caption{The normalised race $E(u)=\frac{\log x}{\sqrt x}(\pi(x;4,3)-\pi(x;4,1))$ for $400\le x\le10^9$ (black), against $u=\log x$, and the model $E_T(u)=1+\sum_{|\gamma|<T}\frac{x^{i\gamma}}\rho$ with the zeros below $T=200$ (red). The dashed line is the bias $1$. Class $1$ leads when $E<0$: near $u=10.2$ ($x=26861$) and a few other places, for $0.04\%$ of the logarithmic scale up to $10^9$.}
\label{fig:race}
\end{figure}

\begin{figure}[ht]
\centering\fig{ch09-density}
\caption{The density of the limiting distribution $\mu$ of Theorem~\ref{thm:rubinsteinsarnak}, computed from its characteristic function with the zeros of $L(s,\chi_{-4})$ below $3000$. It is centred at $1$, with standard deviation $\sqrt V=0.394$. The mass to the left of $0$ is $1-\delta(4;3,1)\approx0.0041$, too small to see at this scale.}
\label{fig:biasdensity}
\end{figure}

\begin{think}
The race modulo $3$, between primes $3n+2$ and $3n+1$, is governed by $L(s,\chi_{-3})$, whose first zero is at height $8.04$, against $6.02$ for $\chi_{-4}$. Before reading Exercise~\ref{ex:mod3}, predict which of the two races is more biased, and why the height of the first zero matters so much.
\end{think}

\section{Other races}

The method applies to every race between two residue classes modulo $q$, with one change: the head start now comes from the number of square roots of each class. Rubinstein and Sarnak proved, under GRH and LI for all the Dirichlet characters modulo $q$, that the logarithmic density $\delta(q;a,b)$ of $\{x:\pi(x;q,a)>\pi(x;q,b)\}$ exists; that it equals $\frac12$ when $a$ and $b$ are both squares or both non-squares modulo $q$; and that it is larger than $\frac12$ when $a$ is a non-square and $b$ a square. For instance $\delta(3;2,1)=0.9990\ldots$ (Exercise~\ref{ex:mod3}). They also proved that $\delta(q;a,b)\to\frac12$ as $q\to\infty$: more and more zeros share the work, and the head start, which comes from a bounded number of square roots, becomes small against their combined oscillation. Daniel Fiorilli and Greg Martin later found asymptotic formulas for these densities and computed many of them to high precision.

The oldest race is also the most famous. The analogue of Proposition~\ref{prop:racel2} for $\zeta$, under the Riemann hypothesis, reads
\[
\frac{\log x}{\sqrt x}\bigl(\pi(x)-\operatorname{li}(x)\bigr)=-1-\sum_\gamma\frac{x^{i\gamma}}\rho+\text{(small in mean square)},
\]
with the same head start of size $1$, coming from the same squares, now on the side of $\operatorname{li}(x)$ (Exercise~\ref{ex:li}). But the zeros of $\zeta$ start higher, at $14.13$, so the waves are weaker, and the bias is crushing: under RH and LI for $\zeta$, the logarithmic density of $\{\pi(x)>\operatorname{li}(x)\}$ is $0.00000026\ldots$ (Rubinstein and Sarnak). No wonder that every table of primes ever computed shows only one side of this race.

\begin{remark}
These are quoted for culture; see the Sources. (a) \emph{Littlewood's theorem} (1914) is unconditional: $\pi(x)-\operatorname{li}(x)$ and $\pi(x;4,1)-\pi(x;4,3)$ change sign infinitely often, and in fact $\pi(x)-\operatorname{li}(x)=\Omega_\pm(\frac{\sqrt x}{\log x}\log\log\log x)$. If GRH fails, the oscillation of size $x^{\beta}$, $\beta>\frac12$, swamps the head start (a theorem of Landau on Dirichlet integrals makes this precise); if GRH holds, Littlewood lines up many waves using Dirichlet's approximation theorem, getting a sum of size $\log\log\log x$, which eventually beats the constant $1$. (b) \emph{Skewes' number.} Littlewood's proof gives no explicit $x$ with $\pi(x)>\operatorname{li}(x)$. His student Stanley Skewes showed in 1933, assuming RH, that there is one below $10^{10^{10^{34}}}$. The bound has since come down by many orders of magnitude: Bays and Hudson (2000) located crossings near $1.398\times10^{316}$, and whether there is an earlier one is still unknown. (c) \emph{Natural density.} The sets in Corollary~\ref{cor:biasdensity} have no natural density: under GRH, Jerzy Kaczorowski showed (see Granville and Martin) that $\frac1X\operatorname{meas}\{x\le X:\pi(x;4,3)>\pi(x;4,1)\}$ oscillates, which is why the logarithmic scale is not a convenience but a necessity (compare Exercise~\ref{ex:logdensity}).
\end{remark}

\begin{beauty}
The bias is an effect of order $\frac{\sqrt x}{\log x}$ in a count of size $\frac x{2\log x}$, invisible to the prime number theorem and to every average over $x$. It survives because of a coincidence of scales: the squares of the primes, weighted as the explicit formula weighs them, give a head start of size $\frac{\sqrt x}{\log x}$, and the waves of the zeros have exactly that amplitude when the generalized Riemann hypothesis holds. The explicit formula turns this coincidence into a probability: the race is a sum of infinitely many independent waves, and the head start is about $2.5$ standard deviations. The same picture, with $\log\zeta$ and the primes, gave Bohr and Jessen's law in Chapter~\ref{ch:bohrjessen}; here the roles are exchanged, and the zeros are the random variables.
\end{beauty}

\section{Exercises}

\begin{exercise}[$\star$]\label{ex:chebbound}
Assume GRH and LI. Using only the mean and variance of $X$ and its symmetry, show that $\delta(4;3,1)\ge1-\frac V2$. With $V=0.1556$, what lower bound do you get?
\end{exercise}

\begin{exercise}[$\star$]\label{ex:logdensity}
Let $A=\bigcup_{k\ge1}[e^{2k},e^{2k+1})$. Show that $A$ has logarithmic density $\frac12$, but that $\frac1X\operatorname{meas}(A\cap[1,X])$ has no limit. Conversely, show that if a set has a natural density, it has the same logarithmic density.
\end{exercise}

\begin{exercise}[$\star\star$]\label{ex:gilpelaez}
Let $Y$ be a real random variable with a symmetric law, no atoms, and characteristic function $\varphi$ with $\int_1^\infty\frac{|\varphi(\xi)|}\xi\dd\xi<\infty$. Show that for $c>0$, $P(Y>-c)=\frac12+\frac1\pi\int_0^\infty\frac{\sin c\xi}\xi\varphi(\xi)\dd\xi$. Show that $\mathbb Ee^{i\xi\cdot2\cos\Theta/r}=J_0(\frac{2\xi}r)$, and deduce the formula for $\delta(4;3,1)$ in the text, assuming the integrability.
\end{exercise}

\begin{exercise}[$\star\star\star$]\label{ex:variance}
Assume GRH. Using Hadamard's factorisation of the entire function $\xi$ of order $1$ (proved as in Tome~5, Theorem~6.7), $\frac{\xi'}\xi(s)=B+\sum_\rho(\frac1{s-\rho}+\frac1\rho)$, show that $\sum_\rho\operatorname{Re}\frac1\rho=-B$ and $B=-\frac{\xi'}\xi(1)$, and deduce
\[
V=\sum_\rho\frac1{|\rho|^2}=\log\frac4\pi-\gamma+2\frac{L'}L(1,\chi_{-4}).
\]
(Numerically, $L(1)=\frac\pi4$ and $L'(1)=0.19290\ldots$, so $V=0.15557\ldots$)
\end{exercise}

\begin{exercise}[$\star\star$]\label{ex:mod3}
Let $\chi_{-3}$ be the non-principal character modulo $3$. (a) Show that $\theta(x;3,2)-\theta(x;3,1)=\theta(\sqrt x)-\psi(x,\chi_{-3})+O(x^{1/3}\log x)$. (b) Explain which changes make the results of this chapter hold for the race modulo $3$, under GRH and LI for $\chi_{-3}$. (c) The analogue of Exercise~\ref{ex:variance} gives $V_3=\log\frac3\pi-\gamma+2\frac{L'}L(1,\chi_{-3})=0.1132\ldots$. Compare the bounds of Exercise~\ref{ex:chebbound} for the two races, and relate the answer to the Think box.
\end{exercise}

\begin{exercise}[$\star\star$]\label{ex:li}
Assume the Riemann hypothesis. Using Tome~5, Theorem~7.5 in place of Theorem~\ref{thm:explicitchi}, show that $\frac u{e^{u/2}}(\pi(e^u)-\operatorname{li}(e^u))=-1-\sum_{|\gamma|<T}\frac{e^{i\gamma u}}\rho+R_T(u)$, with $\frac1U\int_2^U|R_T|^2\to0$ as first $U\to\infty$, then $T\to\infty$. Where does the $-1$ come from?
\end{exercise}

\section{Hints}

\begin{hint}{ex:chebbound}
$P(X\le0)=\frac12P(|X-1|\ge1)$.
\end{hint}
\begin{hint}{ex:logdensity}
In the variable $u=\log x$, $A$ is a union of intervals $[2k,2k+1)$. For the natural density, compare $X=e^{2k}$ and $X=e^{2k+1}$.
\end{hint}
\begin{hint}{ex:gilpelaez}
Write $\mathbf 1\{y>-c\}-\frac12=\frac12\operatorname{sgn}(y+c)$ off $\{y=-c\}$, use $\operatorname{sgn}(z)=\frac2\pi\lim_{M\to\infty}\int_0^M\frac{\sin z\xi}\xi\dd\xi$ with bounded partial integrals, and symmetry to kill the cosine terms.
\end{hint}
\begin{hint}{ex:variance}
Pair $\rho$ with $1-\overline\rho=\overline\rho$ (on the line). The functional equation $\xi(s)=\xi(1-s)$ gives $\frac{\xi'}\xi(0)=-\frac{\xi'}\xi(1)$, and $\frac{\Gamma'}\Gamma(1)=-\gamma$.
\end{hint}
\begin{hint}{ex:mod3}
(a) Every square of a prime $p\ne3$ is $\equiv1\pmod3$. (b) $\chi_{-3}$ is primitive and odd, with root number $1$; the Gamma factor is the same, and only $\frac4\pi$ becomes $\frac3\pi$.
\end{hint}
\begin{hint}{ex:li}
Compare $\pi(x)$ with $J(x)=\sum_{p^k\le x}\frac1k$, and $J$ with $\psi$ by partial summation; the prime squares now appear with a minus sign.
\end{hint}

\section{Solutions}

\begin{solution}{ex:chebbound}
\idea A symmetric variable with small variance rarely strays by its full head start.
\writeup By the proof of Corollary~\ref{cor:biasdensity}, $1-\delta=P(X\le0)=\frac12P(|X-1|\ge1)\le\frac12\mathbb E(X-1)^2=\frac V2$. With $V=0.1556$, $\delta\ge0.922$. The true value $0.9959$ is much closer to $1$ because $X$ is far from the worst case of Chebyshev's inequality: its tails are lighter than Gaussian (each term is bounded).
\end{solution}

\begin{solution}{ex:logdensity}
\idea Logarithmic density is natural density in the variable $u=\log x$.
\writeup In $u$, the set is $\bigcup[2k,2k+1)$, and $\frac1U\operatorname{meas}([0,U]\cap\bigcup[2k,2k+1))\to\frac12$. For the natural density, $\operatorname{meas}(A\cap[1,X])=\sum_{k:e^{2k+1}\le X}(e^{2k+1}-e^{2k})+\dots$; at $X=e^{2K+1}$ this is $(e-1)\sum_{k\le K}e^{2k}\sim\frac{(e-1)e^{2K}}{1-e^{-2}}$, so the ratio tends to $\frac{e-1}{e(1-e^{-2})}=\frac e{e+1}\approx0.73$; at $X=e^{2K}$, the ratio tends to $\frac{(e-1)e^{-2}}{1-e^{-2}}=\frac1{e+1}\approx0.27$. Conversely, if $m(X)=\operatorname{meas}(A\cap[1,X])=\delta X+o(X)$, then $\int_1^X\mathbf 1_A\frac{\dd x}x=\frac{m(X)}X+\int_1^X\frac{m(x)}{x^2}\dd x=\delta\log X+o(\log X)$.
\end{solution}

\begin{solution}{ex:gilpelaez}
\idea The sign function is a Fourier integral, and symmetry keeps only the sine part.
\writeup For real $z$, $\int_0^M\frac{\sin z\xi}\xi\dd\xi=\int_0^{Mz}\frac{\sin v}v\dd v$ is bounded uniformly in $M$ and $z$ and tends to $\frac\pi2\operatorname{sgn}z$. Since $P(Y=-c)=0$, dominated convergence gives
\[
P(Y>-c)-\frac12=\frac12\mathbb E\operatorname{sgn}(Y+c)=\frac1\pi\lim_M\int_0^M\frac{\mathbb E\sin((Y+c)\xi)}\xi\dd\xi,
\]
after Fubini on $[0,M]$ (the integrand is bounded by $|Y+c|$ near $0$ and by $\frac1\xi$ elsewhere; truncate $|Y|\le K$ first if $Y$ is not integrable, and let $K\to\infty$). By symmetry $\mathbb E\sin(Y\xi)=0$ and $\mathbb E\cos(Y\xi)=\varphi(\xi)$, so $\mathbb E\sin((Y+c)\xi)=\sin(c\xi)\varphi(\xi)$, and the integrability assumption lets $M\to\infty$. For the second part, $\mathbb Ee^{i\xi2\cos\Theta/r}=\frac1{2\pi}\int_0^{2\pi}e^{i\frac{2\xi}r\cos\theta}\dd\theta=J_0(\frac{2\xi}r)$, and independence turns the characteristic function of $Y=X-1$ into the product. Take $c=1$.
\end{solution}

\begin{solution}{ex:variance}
\idea The sum over zeros of $\frac1{|\rho|^2}$ is a value of the logarithmic derivative of $\xi$.
\writeup Under GRH, $1-\rho=\overline\rho$ and $\frac1{|\rho|^2}=\frac1{\rho(1-\rho)}=\frac1\rho+\frac1{1-\rho}=\frac1\rho+\frac1{\overline\rho}=2\operatorname{Re}\frac1\rho$. Put $s=1$ and $s=0$ in Hadamard's formula. At $s=0$, $\frac{\xi'}\xi(0)=B$. At $s=1$, $\frac{\xi'}\xi(1)=B+\sum_\rho(\frac1{1-\rho}+\frac1\rho)$, the series converging absolutely. The functional equation gives $\frac{\xi'}\xi(1)=-\frac{\xi'}\xi(0)=-B$, so $\sum_\rho(\frac1\rho+\frac1{1-\rho})=-2B$, that is $V=-2B=2\frac{\xi'}\xi(1)$. Finally $\frac{\xi'}\xi(s)=\frac12\log\frac4\pi+\frac12\frac{\Gamma'}\Gamma(\frac{s+1}2)+\frac{L'}L(s)$, and at $s=1$, $\frac{\Gamma'}\Gamma(1)=-\gamma$. Hence $V=\log\frac4\pi-\gamma+2\frac{L'}L(1)$. Numerically $0.24156-0.57722+2\cdot0.24561=0.15557$.
\end{solution}

\begin{solution}{ex:mod3}
\idea The smaller the variance, the bigger the bias, and the first zero carries a large share of the variance.
\writeup (a) As in Tome~5, Proposition~7.11: the terms $k=1$ of $\psi(x,\chi_{-3})$ give $\theta(x;3,1)-\theta(x;3,2)$; the terms $k=2$ give $\sum_{p\le\sqrt x,\,p\ne3}\log p=\theta(\sqrt x)-\log3$; the rest is $O(x^{1/3}\log x)$. (b) In Theorem~\ref{thm:explicitchi}, $\xi(s)=(\frac3\pi)^{s/2}\Gamma(\frac{s+1}2)L(s,\chi_{-3})$, with the same trivial zeros; the partial summation bound $|S(u)|\le1$ still holds; everything else is unchanged. (c) The bounds are $\delta\ge1-\frac{0.1556}2=0.922$ and $\delta\ge1-\frac{0.1132}2=0.943$. The term of the first zero contributes $\frac2{|\rho_1|^2}=\frac2{\frac14+6.02^2}=0.055$ to $V$ for $\chi_{-4}$, and $\frac2{\frac14+8.04^2}=0.031$ for $\chi_{-3}$. Low zeros make big waves; a higher first zero means a smaller variance, so the fixed head start $1$ counts for more. The true densities are $0.9959$ and $0.9990$.
\end{solution}

\begin{solution}{ex:li}
\idea The same squares, now on the other side.
\writeup Let $J(x)=\sum_{p^k\le x}\frac1k=\pi(x)+\frac12\pi(\sqrt x)+O(x^{1/3})$. Partial summation from $\psi$ gives $J(x)=\int_2^x\frac{\dd\psi(t)}{\log t}+O(1)$. Tome~5, Theorem~7.5 with $T=x^2$ and the argument of Proposition~\ref{prop:racel2} (with Lemma~\ref{lem:zerowaves} for the zeros of $\zeta$, whose proof uses only the local count of zeros, Tome~5, Corollary~6.10) give $J(e^u)-\operatorname{li}(e^u)=\frac{e^{u/2}}u\bigl(-\sum_{|\gamma|<T}\frac{e^{i\gamma u}}\rho+R(u)\bigr)$, with $R$ small in mean square; here the main term $x$ of $\psi$ produces $\operatorname{li}(x)$, and the conversion from $\psi$ to $J$ costs a factor $1+O(\frac1u)$ as before. Then $\frac12\pi(\sqrt x)=\frac{\sqrt x}{\log x}(1+o(1))$ by the prime number theorem. So $\pi(x)-\operatorname{li}(x)=\frac{\sqrt x}{\log x}(-1-\sum\frac{x^{i\gamma}}\rho+\dots)$: the $-1$ is the half-weighted squares of primes, which $J$ counts and $\pi$ does not.
\end{solution}

\section*{Sources and further reading}

\begin{itemize}
\item \emph{Chebyshev's letter.} P.~L.~Chebyshev, letter to P.~N.~Fuss (1853), published in \emph{Bull.\ Classe Phys.-Math.\ Acad.\ Imp.\ Sci.\ St.\ Pétersbourg} 11 (1853), 208; described in Rubinstein and Sarnak (below) and in Granville and Martin (below). J.~Leech, ``Note on the distribution of prime numbers'', \emph{J.\ London Math.\ Soc.}\ 32 (1957), 56--58.
\item \emph{Littlewood, Hardy, Landau.} J.~E.~Littlewood, ``Sur la distribution des nombres premiers'', \emph{C.\ R.\ Acad.\ Sci.\ Paris} 158 (1914), 1869--1872. G.~H.~Hardy and J.~E.~Littlewood, ``Contributions to the theory of the Riemann zeta-function and the theory of the distribution of primes'', \emph{Acta Math.}\ 41 (1918), 119--196. E.~Landau, ``Über einige ältere Vermutungen und Behauptungen in der Primzahltheorie'', \emph{Math.\ Z.}\ 1 (1918), 1--24 and 213--219. A complete proof of Littlewood's theorem is in A.~E.~Ingham, \emph{The Distribution of Prime Numbers}, Cambridge University Press, 1932, Chapter~V, and in H.~L.~Montgomery and R.~C.~Vaughan, \emph{Multiplicative Number Theory~I}, Cambridge University Press, 2007, Chapter~15.
\item \emph{Skewes and after.} S.~Skewes, ``On the difference $\pi(x)-\operatorname{li}x$ (I)'', \emph{J.\ London Math.\ Soc.}\ 8 (1933), 277--283. C.~Bays and R.~H.~Hudson, ``A new bound for the smallest $x$ with $\pi(x)>\operatorname{li}(x)$'', \emph{Math.\ Comp.}\ 69 (2000), 1285--1296.
\item \emph{The limiting distribution.} A.~Wintner, ``On the asymptotic distribution of the remainder term of the prime-number theorem'', \emph{Amer.\ J.\ Math.}\ 57 (1935), 534--538. M.~Rubinstein and P.~Sarnak, ``Chebyshev's bias'', \emph{Experiment.\ Math.}\ 3 (1994), 173--197, whose argument this chapter follows.
\item \emph{Surveys and refinements.} A.~Granville and G.~Martin, ``Prime number races'', \emph{Amer.\ Math.\ Monthly} 113 (2006), 1--33, a delightful introduction with the history. D.~Fiorilli and G.~Martin, ``Inequities in the Shanks--Rényi prime number race: an asymptotic formula for the densities'', \emph{J.\ Reine Angew.\ Math.}\ 676 (2013), 121--212.
\end{itemize}
